\section{Проверка гипотезы о равенстве математических ожиданий \(X_1\) и \(X_2\)}

\subsection{Постановка задачи}

Проверяется гипотеза о равенстве математических ожиданий двух независимых выборок \(X_1\) и \(X_2\).

\subsubsection{Гипотезы}
\begin{itemize}
    \item \(H_0: \mathbb{E}X_1 = \mathbb{E}X_2\) — математические ожидания равны.
    \item \(H_1: \mathbb{E}X_1 \neq \mathbb{E}X_2\) — математические ожидания различаются (двусторонняя альтернатива).
\end{itemize}

\subsubsection{Уровень значимости}
\(\alpha = 0.05\).

\subsection{Выбор критерия и обоснование}

\textbf{Выбранный критерий:} двухвыборочный t-критерий Стьюдента для независимых выборок с равными неизвестными дисперсиями (приложение A.4), реализованный функцией \texttt{scipy.stats.ttest\_ind} с параметром \texttt{equal\_var=True}.

Обоснование применимости:
\begin{itemize}
    \item Выборочные дисперсии близки (\(\hat{\sigma}_1^2 = 74.33\) и \(\hat{\sigma}_2^2 = 86.49\)), что позволяет предполагать равенство дисперсий и использовать pooled-оценку, как в приложении A.4.
    \item Объём выборок достаточно велик (\(n_1 = n_2 = 106\)); согласно центральной предельной теореме, распределение выборочных средних близко к нормальному, что обеспечивает устойчивость t-критерия к умеренным отклонениям от нормальности.
\end{itemize}

\subsection{Вычисление статистики критерия}

Согласно приложению A.4, обозначим выборочные средние и объединённую дисперсию:
\[
\bar{X}_1 = 67.200, \quad \bar{X}_2 = 69.102,
\]
\[
\hat{\sigma}^2 = \frac{(n_1 - 1)\hat{\sigma}_1^2 + (n_2 - 1)\hat{\sigma}_2^2}{n_1 + n_2 - 2}
= \frac{105 \cdot 74.329 + 105 \cdot 86.489}{210}
\approx 80.409.
\]

Статистика критерия (приложение A.4):
\[
h = \frac{\bar{X}_1 - \bar{X}_2}{\hat{\sigma} \sqrt{\frac{1}{n_1} + \frac{1}{n_2}}}
= \frac{67.200 - 69.102}{\sqrt{80.409} \cdot \sqrt{\frac{2}{106}}}
\approx -1.544.
\]

При справедливости \(H_0\) статистика \(h\) имеет распределение Стьюдента с \(\nu = n_1 + n_2 - 2 = 210\) степенями свободы.

\subsection{Принятие решения}

Согласно приложению A.3, используем p-value:
\[
p\text{-value} = 2 \cdot P(T_{210} \leq -1.544) \approx 0.124.
\]

\[
p\text{-value} = 0.124 > \alpha = 0.05.
\]

Критическая область для двусторонней альтернативы (приложение A.4):
\[
W = (-\infty, -t_{210,\, 1-\alpha/2}] \cup [t_{210,\, 1-\alpha/2}, +\infty)
= (-\infty, -1.971] \cup [1.971, +\infty).
\]

Наблюдаемое значение \(|h| = 1.544 < 1.971 = t_{210,\, 0.975}\), следовательно \(h \notin W\) — нет оснований отвергнуть \(H_0\). Оба подхода (p-value и критическая область) дают одинаковый результат.

\subsection{Вывод}

\textbf{Нет оснований отвергнуть нулевую гипотезу} \(H_0: \mathbb{E}X_1 = \mathbb{E}X_2\) на уровне значимости \(\alpha = 0.05\). Статистически значимых различий между средними \(X_1\) и \(X_2\) не обнаружено.

\subsection{Ошибка первого рода}

В данной задаче ошибка первого рода (приложение A.2) означает: \textbf{отвергнуть гипотезу о равенстве математических ожиданий, когда на самом деле они равны}. Вероятность такой ошибки контролируется уровнем значимости \(\alpha = 0.05\).
